% 本文件由 LaTeX 排版 Agent 根据有效 translation.json 自动生成。
% author=Cyprian of Carthage; work_id=0505; title=圣西彼廉的生平与受难

\chapter{圣西彼廉的生平与受难}

% structure: p0001 source=h1 target=omitted reason=page-title-already-rendered-by-parent

\Signature{执事彭提乌斯}\par

1. 虽然西彼廉，这位虔诚的司铎和光荣的天主见证人，撰写了众多著作，使其令名之记忆得以流传；虽然他的口才与天主恩宠的丰饶多产，以其论述的丰沛与富足如此扩展，以至于他可能直至世界终了都不会停止说话；然而，既然他的事迹与功绩正当地要求将其榜样以文字记录下来，我认为准备这份简短而扼要的叙述是合适的。并非如此伟大之人的生平甚至对外邦民族中的任何人都是未知的，而是为了使我们的后代也能将这无与伦比、崇高的典范延续到不朽的记忆中。这肯定是一件艰难的事，当我们先辈们甚至对获得殉道恩典的平信徒和慕道者也给予了如此尊荣，出于对他们殉道的敬畏，以至于记录了许多——我几乎可以说，几乎全部——他们受苦的细节，以便它们也能为我们这些当时尚未出生的人所知，那么像西彼廉这样的司铎和殉道者的受难竟被忽略：他，独立于他的殉道，还有许多东西要教导，并且他生前所行之事不应向世界隐藏。并且，的确，他的这些事迹如此伟大、如此令人钦佩，以至于我因思虑其伟大而被震慑，承认自己无力以配得其功绩尊荣的方式论述，也无法以如此方式叙述这些高尚事迹，使它们显得与事实一样伟大，除非他的荣耀众多，本身已足够，无需其他传扬。这增加了我的困难，因为你也渴望听到关于他的许多事情，如果可能，每一件事，热切地渴望至少了解他的事迹，尽管如今他活生生的言语已然沉默。在这方面，如果我说我的口才之力不足，我说的还太少。因为口才本身也缺乏足够的力量来完全满足你的愿望。因此我两面受困，因为他以他的美德重压我，而你以你的恳求紧逼我。\par

2. 那么，我从何处开始？从哪个方向着手描述他的美善，除非从他的信德之始和从天上的诞生开始？因为天主之人的事迹不应从任何时刻算起，除了从他由天主所生的那一刻。他之前可能有过追求，人文学科可能在他从事其中时浸润了他的心灵；但这些东西我略过不提；因为当时它们除了他的世俗利益外与任何事无关。但当他学习了圣学，冲破此世云雾而进入属灵智慧之光时，如果我曾与他一同参与他的任何事迹，如果我曾辨识出他任何更杰出的劳苦，我将述说它们；同时只请求这一宽容：凡我说得不足的（因为必然说得不足），与其归咎于我的无知，不如归为不减损他的荣耀。当他的信德尚在初阶时，他相信在天主面前，没有什么比遵守法律更值得。因为他认为，只有以圣德的强健活力践踏肉身的私欲，心灵才能成为它应有的样子，心智才能达到真理的圆满能力。谁曾记录这样的奇事？他的重生尚未以神圣之光的全部光辉启迪这新人，但他已仅凭曙光就战胜了古老原始之黑暗。然后——甚至更伟大的是——当他从阅读圣经中学到某些事，不是按他初学阶段的条件，而是按他信德的早熟程度，他立刻抓住他所发现的，为了自己堪得天主青睐的益处。通过分配他的资产以救济穷人的匮乏，通过变卖全部田产所得的款项，他立即实现了两个益处——轻看此世的野心，没有比这更危险的，以及遵守那怜悯，即天主甚至优先于祭献的怜悯，而且连那说自己遵守了法律一切诫命的人也没有持守这怜悯；由此，以虔诚的过早迅速，他几乎在他学会如何成为完美之前就开始变得完美。请问，古人中谁曾这样做？在信德中最著名的老兵中，他们的心与耳多年来为神圣的话语而悸动，有谁曾尝试过这样的事，如同这人——信德尚未熟练，也许尚无人信任——超越古代的年岁，以他光荣而令人钦佩的劳苦所成就的？没有人播种后立即收割；没有人从刚挖的沟中压榨出葡萄收获；从未有人从新插的枝条中寻求成熟的果实。但在他身上，所有难以置信之事同时发生。在他身上，打谷（如果可以这样说，因为这难以置信）——先于播种，葡萄收获先于嫩芽，果实先于根。\par

3. 宗徒的书信说\BibleRef{弟前 3:6}，新手应被忽略，以免仍附着于他们未经坚固的心灵的异教愚昧，使他们的未经教导的经验在任何事上得罪天主。他首先，且我认为只有他，提供了一个例证：信德比时间取得更大的进步。因为虽然在\BibleBook{宗徒}{大事录}中，太监被描述为由\Person{斐理伯}立刻受洗，因为他全心相信，但这并非一个公平的类比。因为他是犹太人，且从主的圣殿而来，他正在阅读\BibleBook{依撒意亚}{先知书}，并希望于基督，尽管他尚未相信基督已来临；而另一个来自无知的异教徒，开始时拥有一种成熟的信德，这种信德或许很少有人能坚持到底。简言之，关于天主的恩宠，没有迟延，没有拖延——我说得还不够——他立刻接受了长老职与司铎职。因为谁不会将各种荣誉托付给一个怀着如此心志相信的人呢？他还在平信徒时就做了许多事，现在作为长老也做了许多事——按照古时义人的榜样，并紧密效法他们，他以完全奉献的服从完成了这些事——这些事在主面前应得奖赏。因为关于这一点，他通常这样谈论：如果他读到任何一个人被天主赞扬，他会劝我们探究那人因何行为而中悦天主。如果\Person{约伯}，由天主的见证而光荣，被称为天主的真敬拜者，且世上无人能与他相比，他就教导我们应做\Person{约伯}以前所做的一切，这样当我们做类似的事时，也能为自己求得天主类似的见证。他蔑视产业的损失，因这被考验的德行而获益如此之多，以至于他甚至没有感受到情感上暂时失落的痛苦。贫穷和痛苦都未能击倒他；妻子的劝说没有影响他；自己身体可怕的痛苦没有动摇他的坚定。他的德行在自身的根基中稳固，他的虔诚建立在深根之上，在魔鬼诱惑他即使在逆境中也停止以感恩的信德赞美天主的任何攻击下毫不退缩。他的房屋对每一位来者开放。没有一个寡妇空着衣襟从他那里离开；没有盲人不被他作为同伴引导；没有脚步蹒跚的人不被他作为手杖扶持；没有因强者的手而失去帮助的人不被他作为保护者保护。他惯常说，凡愿意中悦天主的人都应该做这样的事。如此，他通过效法所有善人的榜样，总是模仿那些比别人更好的人，也使自己堪受效法。\par

4. 他与一位义人，一位值得赞美的记忆，名叫\Person{凯齐略}的人，在我们中间有密切的关系，他在年龄和尊荣上都是一位司铎，此人曾将他从世俗的错误中引向对真天主的认识。他以完全的尊敬和服从爱他，以顺从的敬畏看待他，不仅将他视为灵魂的朋友和同伴，而且视为他新生命的父亲。最后，他受他关怀的影响，被激励到如此过度的爱，以至于当他离开此世、他的召唤临近时，他将自己的妻子和儿女托付给他；这样，他既已使他成为自己生活方式的伙伴，后来便使他成为自己感情的继承人。\par

5. 逐一叙述个别情况会令人厌倦，详述他的一切作为也是费力的。为了证明他的善行，我认为这一件事就足够了：由天主的判断和民众的喜爱，他还在新信徒和被认为新手的时候，就被选任为司铎职和主教职。虽然还在信德的早期，以及灵性生命的未经教导的时节，一种慷慨的性情在他身上如此闪耀，以至于尽管尚未焕发职务的光彩，而只是希望，他就为即将到来的司铎职提供了完全可信的保证。此外，我不会忽略那个显著的事实：当全体民众因天主的感动，以爱和尊荣跳跃前进时，他谦卑地退后，将位置让给更年长的人，认为自己不配要求如此大的尊荣，因此他变得更加配得。因为一个拒绝其所应得的人变得更加配得。当时热切的人们被这种兴奋所点燃，渴望见到一个主教——因为在他身上，他们以某种神圣的预兆如此要求，他们不仅是在寻找一位司铎，而且更是一位未来的殉道者。拥挤的弟兄团体包围了房屋的门，焦急的爱在所有通道上流动。也许那宗徒的经验那时会发生在他身上，正如他所愿，从窗户被放下来，如果他在领受祝圣的尊荣上也等同于宗徒的话。可以看到，其他所有人都以焦虑的悬念之精神等待着他的到来，当他来到时，以过度的喜悦接待了他。我不情愿地说，但必须说。有些人反对他，甚至使他胜过他们；然而他是何等温柔，何等耐心，何等仁慈地宽容他们！他是何等仁慈地宽恕他们，后来将他们列为他最亲近、最亲密的朋友，令许多人惊讶！因为谁能为一个如此强记的头脑的善忘而惊讶呢？\par

6. 从此以后，谁能充分叙述他如何自持？——他的怜悯何等深切？他的精力何等充沛？他的仁慈何等伟大？他的严厉何等深厚？从他脸上发出如此多的圣洁和恩宠，使观看者心神错乱。他的面容庄重而喜乐。他的严厉并不阴沉，他的和蔼也不过分，而是两者的和谐调合；以至于令人怀疑，他更应受到敬畏还是爱戴，但他却同时值得敬畏和爱戴。他的衣着也与他的面容相协调，也同样适中。世俗的骄傲没有燃起他，过分的故作贫寒也没有使他卑劣，因为后一种服饰同样出自夸耀，正如那种野心勃勃的节俭出自炫耀。但作为主教，他对待那些他曾作为望教者所爱的穷人，做了些什么呢？让虔诚的司铎们思考，或那些因其职位的教导而被培养行善之责的人，或那些因圣事的共同义务而被约束以显明爱德的人。西彼廉主教的\OriginalTerm{主教座位}{cathedra}接纳了如他先前那样的人——并未使他成为那样。\par

7. 因此，由于这样的功德，他立即也获得了被放逐的光荣。因为最恰当的是，那在良心隐秘处富有宗教和信德之全部尊荣的人，还应在外邦人公开传播的言论中闻名。确实，按照他恒常迅速达成一切的速度，他本可在那时就奔向为他预定的殉道荣冠，尤其当他屡次被呼、常被要求投入狮口时；但若非他必须经过所有光荣的等级，从而达到最高者，且若非即将来临的荒凉需要如此丰饶心智的帮助，则他早已如此。请设想他当时被殉道的尊威带走：谁将显示由信德而前进的恩宠之益处？谁将约束贞女们持守适度的端庄和与圣洁相称的衣饰，仿佛用主之教训的缰绳？谁将教导背教者补赎、异端者真理、分裂者合一、天主子女们平安与福音祈祷的规范？谁将以他们加于我们的控告反击自身，从而战胜亵渎的外邦人？谁将以未来的希望安慰那些对亲人丧失感情过于柔嫩、或更重要的是信德过于软弱的基督徒？我们从何处如此学习仁慈？从何处学习忍耐？谁将以有益良药的甘甜遏制由嫉妒的恶毒所产生的坏血？谁将以神圣言论的劝勉兴起如此伟大的殉道者？简言之，谁将激励那么多以第二道铭文封于卓越额头、并被保留活着作为殉道榜样的宣信者，以天国的号角点燃他们的热忱？幸运，幸运地那时发生了，并且确实是藉圣神的引领，那为如此众多且卓越目的所需之人被阻止了殉道的完成。你愿确信这退避的原因不是恐惧吗？别的姑且不提，他后来确实受苦了，而他若先前曾逃避，无疑也会如常逃避这次受苦。的确，那是恐惧——而且合理——那种畏惧得罪主的恐惧，那种宁愿服从天主的命令而不愿在不服从中被加冕的恐惧。因为那全心奉献于天主并如此服从于神圣训诫的心灵，相信即使在受苦本身中也会犯罪，除非它服从了主，那时主命令\Emph{他寻找}藏身之处。\par

8. 此外，我认为此处可以略述延迟的益处，尽管我已稍微触及此事。从后来发生的事看来，我们可以证明那退避并非出于人性的怯懦，而是——诚然如此——真正神圣的。一场残酷迫害的不寻常猛烈愤怒蹂躏了天主的子民；由于狡猾的敌人不能以一诈欺欺骗所有人，无论粗心的战士在哪里暴露他的侧翼，那里就以各种愤怒的表现，用不同的毁坏方式毁掉个人。需要一个人，能在人们被攻击之敌的各种伎俩伤害时，根据伤口性质使用天国药物的疗法，或切割或抚慰他们。\Emph{如此}保存了一个智慧超群并受过属灵训练的，在反对分裂的涛声之间能将教会稳妥地航行于中道的人。我问，这些计划岂非神圣？没有天主能成就这事吗？让那些以为此类事情出于偶然的人思考。教会以清晰的声音回答他们说：\Quote{我不允许也不相信，这样的需要那时竟没有天主的旨意而得以保存。}\par

9. 不过，若合适，让我简单看一下其余的事。随后爆发了一场可怕的瘟疫，一种可憎疾病的过度毁灭接连侵入战栗的人群的每一个家庭，每天以突然的攻击带走无数人，人人从自己的家中被带走。所有人都战栗、逃窜、躲避传染，不虔地暴露自己的朋友，仿佛通过排除注定死于瘟疫的人，也能排除死亡本身。与此同时，整个城市里躺着的不再是身体，而是许多人的尸骸，并通过对自己也将遭遇的命运的默想，向过路人求取怜悯。没有人顾及任何事，除了他残忍的利益。没有人因对类似事件的回忆而战栗。没有人对他人做他自己希望经历的事。在这种情况下，忽略\OriginalTerm{基督的主教}{pontiff of Christ}所做的，是不对的；他在仁慈的关爱和信仰的真理性上，都胜过\OriginalTerm{世界的主教}{pontiffs of the world}。他首先向聚集在一处的民众敦促仁慈的益处，通过神圣读经中的例子教训他们，仁爱的职分对于赢得天主的好感有极大助益。随后他又补充说，我们只以必要的爱的关怀照顾我们自己的人，并没什么了不起；但若有人做得比税吏或外邦人更多，以善胜恶，并实践一种类似天主仁慈的仁慈，甚至爱他的仇敌，为主所劝诫和鼓励的，为迫害他的人祈祷得救，那么他就可能成为成全的人。天主不断地使太阳升起，并时常降下雨水滋润种子，显示所有这些恩惠不仅给祂的子民，也给外邦人。若有人自称是天主的儿子，为什么他不效法他父亲的榜样呢？\Quote{适合我们的，}他说，\Quote{是回应我们的出生；那些显然由天主所生的人不应堕落，而应藉着对祂美善的效仿，在祂的后代中证明一位良善圣父的繁衍。}\par

10. 我省略了许多其他事情，确实许多重要的事情，有限篇幅的必要不允许用更长的论述详述它们，而说到这些就足够了。但如果外邦人当时站在讲台前能听到这些，他们很可能立刻就会相信。那么，一个基督的百姓应当做什么呢？他们的名字本身就来自信德。于是，服事常常根据人的品级和程度分配。许多因贫穷的窘迫而无法显示财富之仁慈的人，显示了比财富更多的东西，用自己的辛劳弥补了一种比一切财富更珍贵的服务。在这样一位导师的带领下，谁不想在这样一场战争中的某个部分被找到，从而可以同时中悦天主圣父、基督审判者，并为了现时如此卓越的一位司铎呢？因此，良善以洋溢工作的慷慨行于众人，不仅行于信德之家的人。所做的比关于\Person{多俾亚}无与伦比的慈惠的记载更多。他必须宽恕，再宽恕，常常宽恕；或者更准确地说，他当然必须承认，虽然许多事可以在基督之前做，但更多的事可以在基督之后做，因为祂的时代拥有一切圆满。\Person{多俾亚}只聚集了那些被国王杀害并被抛弃的人，都是他本族的。\par

11. 在如此良善与仁慈的行动之后，随之而来的是放逐。因为不虔敬总是以这样的方式回报：以更恶报更善。至于天主的司铎对总督审讯所作的回答，有\WorkTitle{行实}记载。与此同时，那位曾为城市的安全做了些好事的人被逐出城外；他曾努力使活人的眼睛不致遭受阴间住所的恐怖；我说，他曾在仁爱的警醒中看守，以不被承认的善行——何其邪恶！——在所有人都遗弃城市荒凉的外貌时，使一个贫困的城邦和一片被遗弃的土地不致察觉其众多的流放者。但让世界去审视这一点，它视放逐为刑罚。对他们而言，祖国极为珍贵，他们与父母同姓；但我们甚至憎恶我们的父母，如果他们想说服我们反对天主。对他们而言，住在自己的城邦之外是严厉的惩罚；对基督徒而言，整个世界就是一个家。因此，即使他被放逐到一个隐蔽隐秘的地方，既与他的天主的事相连，他便不能视之为流放。此外，在诚心事奉天主的同时，即使在自己的城中他也是个陌生人。因为圣神的节制使他远离肉体的欲望，他放下了旧人的交谈，即使在同乡中间，或者我几乎可以说，在他世俗生活的父母中间，他也是个陌生人。此外，尽管这在其他方面看似一种惩罚，但在此类我们为验证自身美德而忍受的缘由和判决中，它并非惩罚，因为是荣耀。但确实，假设放逐对我们不是惩罚，他们自己的良心的见证仍可将最坏至极的邪恶归咎于那些能够将他们认为的惩罚加于无辜者的人。我现在不描述一个迷人的地方；暂时略过所有可能的欢乐的增加。让我们想象这地方：位置污秽，外观肮脏，没有有益健康的水，没有悦目的青翠，没有邻近的海岸，而是在一片完全荒凉的荒野的无好客的巨口之间，有广袤的岩石林地，远离世界无路可寻的区域。这样的地方本可承受流放之名，如果天主的司铎\Person{西彼廉}曾来到那里；尽管对他而言，如果缺乏人的服侍，要么如\Person{厄里亚}时的乌鸦，要么如\Person{达尼尔}时的天使，都会服侍他。去，去相信任何事物会缺乏我们中最小的一个，只要他站在为宣认主名的立场上。天主的这位主教，始终急于行仁慈的工作，远非需要所有这些事物的帮助。\par

12. 现在让我们满怀感激地回到我原先提出第二点：为这样一个人的灵魂，天主上智地预备了一处阳光明媚、合适的地方，一所如他所愿的隐秘居所，以及一切先前应许要加给那些寻求天国和天主之义的人的东西。并且，且不提我去探访他的弟兄们的数目，以及市民们本身的善意，他们供应了他一切看似被剥夺的东西，我也不愿略过天主奇妙的眷顾：祂愿意祂的司铎在放逐中如此确定随后将要发生的受难，以致在他对那威胁性殉道的完全信赖中，\Place{库鲁比斯}不仅拥有一个流放者，也拥有一个殉道者。因为在我们初次居住于放逐之地的那一天（因他爱的俯就，他选择了我作为他家中同伴，去自愿的放逐：但愿他也能选择我分享他的受难！），\Quote{有一个少年出现在我面前，}他说，\Quote{那时我尚未沉入睡眠的安息，一个身材异常高大的少年，他仿佛领我到了总督府，在那里我自己似乎被带到当时正坐堂的总督的法庭前。当他注视我时，便立刻开始在记事板上写下一个判决，但我不知道那是什么，因为他没有按惯常的审讯问我什么。但那站在他背后的少年非常焦虑地读着所记下的内容。因他那时不能用言语宣告，他便通过可理解的记号向我指明那记事板上所写的内容。因为他将手展开并压平如刀片，模仿了惯常刑罚的一击，如言语般地清晰表达了他想被理解的事——我明白了未来我受难的判决。我立即开始请求并乞求至少应许我一日的延迟，直到我以合理的方式安排我的财产。当我急切地重复我的恳求时，他又开始在记事板上记下我不知道的东西。但我从他面容的平静觉察到，审判者的心被我的恳求所动，认为它是合理的。此外，那个已经通过手势而不是言语向我透露了我受难消息的少年，急忙通过秘密信号反复示意所请求的延迟已被应允直到次日，他将手指一根一根地扭在一起。而我，虽然判决未被宣读，虽然我因所获的延迟而满心欢喜，却因对解释的不确定而如此恐惧战栗，以致恐惧的残余仍使我狂喜的心因过度激动而跳动。}\par

13. 有比这启示更清楚的么？有什么比这屈尊就卑更蒙福的么？一切都在他预先被告知，而随后都一一应验了。天主的话语没有一句落空，如此神圣的恩许也没有任何减弱。请仔细思量每一细节如何与所宣告的相符。当他受难判决正在讨论时，他请求延期至次日，恳求能在他所获得的这一天内安排自己的事务。这一天象征一年，即他在异象之后将要留在世上的时间。说得更明白些，一年期满后，他正是在当年年初这件事被宣告的那一天获得了殉道的荣冠。因为，尽管我们在圣经中不常将主的日子读作一年，但我们认为这段时期在应许未来之事上是当得的。既然如此，若在此处，在\Quote{一天}的通常说法下，实际暗示的只是一年，又有什么关系呢？因为更重大的事物应当含有更丰满的意义。再者，此事用标记而不是用言语来解释，是因为言语的表达被保留到时间本身显现时。因为，凡所宣告的事物，通常都是在完成之时才用言语说出来。实际上，没有谁知道为何此事向他显示，直到后来，就在他看见异象的同一天，他获得了殉道的荣冠。然而在此期间，他即将面临的苦难当然为众人所知，但没有人说出他受难的确切日子，仿佛他们不知情似的。确实，我在圣经中也发现类似的情况。司祭\Person{匝加利亚}，因为不相信天使向他预告得子的恩许，成了哑巴；以致他用手势要求写字板，准备写下他儿子的名字，而不说出来。有理地，在此情形中，天主的使者用标记而非言语宣告祂司铎即将来临的受难，他既劝勉了他的信德，也坚固了他的司铎。此外，请求延期的缘由出于他想安排事务并立定遗嘱。但他有什么事务或遗嘱要安排，除非是关于教会的事？因此，那最后的延期被允许，是为了使他能以最后的决定安排一切关于照顾穷人需处理的事。我认为没有其他原因，确实仅为这个原因，甚至那些驱逐他并即将杀他的人也给予他宽容，使他能在现场用他最后管家职的最终支出，或更准确地说，全部支出，来救济他面前的穷人。因此，如此慈祥地处理事务，并照他的意愿安排妥当后，翌日临近了。\par

14. 这时也有一位使者从城里来到他这里，是来自\Person{西斯笃}，那位良善且缔造和平的司铎，因此也是最蒙福的殉道者。刽子手立刻被期待前来，要刺穿那最神圣祭品已经奉献的脖颈；如此，在每日期待死亡中，每一天在他看来都像是可能归于他的荣冠。同时，有许多显要人士、出身最显赫的家族、在世间的尊荣中高贵的人聚集到他那里，因着与他旧日的友谊，一再催促他退避；且为使他们的催促不至空洞，他们还提供了他可以退隐的地方。但他现在已将世界置诸脑后，心灵悬于天上，并不应允他们诱惑性的劝说。如果他被天主的命令所吩咐，他甚至可能已经答应了这么多忠实朋友所请求的事。但如此伟大人物的崇高荣耀不可不加宣告就被忽略：当世界正在膨胀，并因信赖其王侯而吐出对圣名的仇恨时，他正借着机会，用主的劝言教导天主的仆人，并激励他们借着默观将来要来的光荣，践踏今时的苦难。的确，他如此热爱神圣的讲论，以致他盼望自己的苦难得以如此应允，即他正在谈论天主时被处死。\par

15. 这些正是一位司铎的日常行为，他注定要成为天主悦纳的祭献。看哪，当总督下令时，军官与士兵突然出乎意料地来到他面前——或者更准确地说，自以为出其不意地来到他面前——来到他的园囿。我是说，来到他的园囿，就是他在信德之初曾变卖、后又因天主的仁慈而复得的园囿；他本来必定会再次变卖以供穷人使用，若不是他想要避免招致迫害者们的恶感。然而，一颗时刻准备好的心神，何时会被突如其来的攻击所惊扰呢？因此，他现在前行，确信那久已延迟的事将要了结。他昂首阔步地前行，面容显露喜乐，心中充满勇气。但因为被推迟到次日，他从总督府返回军官的住所。突然，一个零散的传言传遍了整个\Place{迦太基}，说\Person{Thascius}已被带出。没有人不认识他，既因其辉煌的声名深得众人敬重，也因其极为卓越的工作而被人怀念。四面八方的人们蜂拥而至，要观看这一景象——这景象对我们来说，因信德的虔诚而光荣，甚至对外邦人来说也是可哀叹的。然而，当他被带走并被安置在军官住所过一夜时，一种温和的看管负责看守他，以至于我们这些他的同伴和朋友，如往常一样与他在一起。与此同时，全体人民焦虑不安，唯恐整夜有任何事情在他们不知情的情况下发生，他们守在军官的门口守望。天主的恩慈在那时刻赐予他——他如此地堪当此恩——以至天主的子民竟为司铎的殉道而守望。然而，或许有人会问他为何从总督府返回军官那里。有人认为这是因为总督当时不愿亲自主理。我绝不敢在神圣安排的事上抱怨总督的懈怠或厌恶。我绝不敢让这样一种恶念进入虔诚心灵的意识中，以为人的幻想能决定如此有福的殉道者的命运。但是，早在一年前天主的俯就所预言的次日，必须就是在字面意义上的次日。\par

16. 最后，那另一天破晓了——那命定的、应许的、神圣的一天——即使暴君自己愿意推迟，他也无力为之；这一天因预知未来的殉道者而喜乐；乌云驱散，环绕世界，那天用灿烂的太阳照耀他们。他从军官的住所出来，尽管他本是基督和天主的军官，四周被混杂的人群所围困。如此无数的队伍紧随着他，仿佛集结的军队要来攻击死亡本身。当他前行时，他必须经过赛马场。这是公平的，仿佛是刻意安排的，他必须经过这相称的搏斗场所，因为他已结束竞赛，正奔向正义的冠冕。但他到达总督府时，总督尚未出来，于是赐给他一处休息的地方。在那里，他坐着，长途跋涉后浑身被汗湿透（座位恰好铺着亚麻布，使得就在他殉道的时刻，他也能享有主教职位的尊荣），一位军官（\Quote{Tesserarius}），曾是基督徒，将自己的衣服递给他，似乎想让他把湿衣服换得更干爽；而他无疑除了要获得那即将走向天主的殉道者如今被血沾湿的汗水之外，别无所求。他回答他说：\Quote{我们用药来治疗今日或许不再存在的烦恼。}谁能怪罪他轻视肉体的痛苦呢？他已在灵魂上轻视了死亡。我们何必再多说？他忽然被禀报给总督；他被带上前；他被置于总督面前；他被问及姓名。他回答了他是谁，别无他言。\par

17. 因此，审判官从他的案板宣读判决，这判决是近日在异象中他不曾读过的——一个属神的判决，不可轻率宣告——一个堪当如此主教与如此见证者的判决；一个荣耀的判决，其中他被称作该教派的标准旗手、众神之敌、以及他子民的榜样；并且，藉着他的血，纪律将开始被建立。没有比这判决更完全、更真实的了。因为所说的一切，虽然出自外教人之口，却是神圣的。这并不足为奇，因为司铎惯于预言殉道。他曾是标准旗手，惯于教导如何背负基督的旗帜；他曾是众神之敌，曾命令摧毁偶像。此外，他给他的朋友们作了榜样，因为当许多人将要效法他时，他是在该省首先奉献了殉道的初果。藉着他的血，纪律开始被建立；但那乃是殉道者的纪律，他们效法自己的老师，在模仿他的荣耀中，他们自己也藉着自己榜样的血，给纪律作了确认。\par

18. 当他离开\Place{总督府}的大门时，一群士兵陪同着他；为了使他的苦难毫无欠缺，百夫长和护民官守卫在他身旁。他即将受苦的地方本身是平坦的，因而呈现出一幅壮观的景象，四周密植树木。但是，由于远处空间的广阔，混乱的人群无法看到，那些支持他的人爬上了树枝，这样他也不会缺少（正如\Person{匝凯}身上发生的事）从树上被人注视。现在，他亲手蒙住自己的眼睛，试图催促刽子手的迟缓，刽子手的职责是\Emph{挥舞}刀剑，他颤抖的手指艰难地握住衰老右手中的剑刃，直到光荣的成熟时刻用从上而来的权能坚固了百夫长的手，以完成这位杰出者的死亡，并最终赐予他被允许的力量。哦，教会蒙福的子民，你们不仅用眼睛和感觉，更用直言不讳的言语，与这样一位主教一同受苦；正如他们曾在他自己的讲道中所听到的，由\Person{天主}审判者加冕！因为尽管普遍愿望所渴望的事未能发生，即全体会众应在相似光荣的共融中同时受苦，然而谁若在\Person{基督}的注目下和司铎的听闻中，热切渴望受苦，凭那渴望的充足见证，就在某种程度上像他的使者一样，给\Person{天主}发送了一封信函。\par

19. 他的苦难就这样完成了，结果是\Person{西彼廉}，他曾是所有善人的榜样，也成为在\Place{非洲}第一个以\Emph{殉道者的血}染红了他的司铎冠的人，因为他是宗徒之后第一个开始如此的人。因为自从在\Place{迦太基}的主教职列开始枚举时，从未记载过任何人，即使是善人和司铎，曾达到受苦。虽然奉献给\Person{天主}的虔诚在献身者中常被视作殉道，但\Person{西彼廉}却藉着\Person{主}的完成达到了完美的冠冕；因此，就在他曾如此生活的那座城市，在那里他曾是许多高尚事迹的第一个，他也第一个用光荣的血装饰了他天上司职的徽章。我现在该怎么办？在他的苦难带来的喜悦和仍然存留的悲伤之间，我的心灵被不同方向分裂，双重的感情重压着一颗对它们来说过于有限的心。我该为我不是他的同伴而悲伤吗？但我仍须在他的胜利中欢欣。我该为他的胜利欢欣吗？我仍悲伤我不是他的同伴。然而我仍必须坦诚地向你们告白，你们也知道，那就是我曾打算成为他的同伴。我极为过度地欢欣于他的光荣，但更甚的是，我为未能同行而悲伤。\par

\SourceRecord{Translated by Robert Ernest Wallis.；Ante-Nicene Fathers；Vol. 5.；Edited by Alexander Roberts, James Donaldson, and A. Cleveland Coxe.；Buffalo, NY: Christian Literature Publishing Co.,；1886.；<http://www.newadvent.org/fathers/0505.htm>.}
